\documentclass[11pt,a4paper]{article}
\usepackage[margin=0.92in]{geometry}
\usepackage{amsmath,amssymb}
\usepackage{booktabs}
\usepackage{array}
\usepackage{xcolor}
\usepackage{hyperref}
\usepackage{longtable}
\hypersetup{colorlinks=true,linkcolor=blue!55!black,citecolor=blue!55!black,urlcolor=blue!55!black}
\setlength{\parindent}{2em}
\setlength{\parskip}{0.30em}
\setlength{\emergencystretch}{2em}

\newcommand{\Rzero}{\mathcal R_0}
\newcommand{\Afun}{\mathcal A_\theta}
\newcommand{\Jfun}{\mathcal J_\theta}
\newcommand{\Cfun}{C_\theta}
\newcommand{\Dfun}{D_\theta}
\newcommand{\BL}{y_{\rm BL}}
\newcommand{\xin}{x_{\rm in}}
\newcommand{\xf}{x_f}
\newcommand{\xs}{x_*}
\newcommand{\Pone}{P_1}
\newcommand{\Pburn}{P_{\rm burn}}

\title{Epidemic Burnout under General Susceptible Recruitment\\
\large A Unified Asymptotic Theory from Phase-Plane Matching to Second-Order $x_{\rm in}$ and Kendall Persistence Probability}
\author{Unified theory note}
\date{2026-08-13}

\begin{document}
\maketitle

\begin{abstract}
We develop a self-contained asymptotic framework for the probability of pathogen burnout or persistence after the first epidemic wave. We consider
\[
\frac{dx}{d\sigma}=\rho(1-x)x^\theta-\Rzero xy,
\qquad
\frac{dy}{d\sigma}=(\Rzero x-1)y,
\qquad 0\leq\theta\leq1.
\]
The central problem is to estimate the susceptible fraction $\xin$ when the post-peak infectious fraction first falls to a low-prevalence matching line $y=\BL$, without integrating the full two-dimensional epidemic ODE, and then to use that estimate to calculate the probability $\Pone$ that the pathogen survives the first epidemic trough.

Following the matched-asymptotic phase-plane architecture of Parsons and Earn, we extend standard linear susceptible replenishment to the family $h_\theta(x)=(1-x)x^\theta$. The zero-demography epidemic geometry first gives a universal outer trajectory $F(x)$ and its final-size root $x_f$. We then directly expand the exact phase-plane equation in the fixed-$\Rzero>1$, $\rho\to0$ limit. The first-order correction develops a logarithmic singularity near $x_f$; its finite part defines the general-$\theta$ refined matching amplitude $C_\theta$, which reduces exactly to the refined Parsons et al. burnout prefactor when $\theta=0$. Continuing the same expansion to $O(\rho^2)$ yields a second matching constant $D_\theta$ and the next corner switchback. Re-expressing the corner expansion in terms of the action primitive $\mathcal A_\theta$ gives a stable second-order action-resummed scalar equation for $x_{\rm in}$.

Finally, $x_{\rm in}$ feeds directly into Kendall's time-inhomogeneous birth--death formula. A high-accuracy version requires only a one-dimensional Kendall integral; when the action is sufficiently large, Laplace's method gives a fully asymptotic closed form. We also distinguish several boundaries of validity: no boundary-layer entry, scale merging as $\Rzero\to1^+$, small-action Laplace failure, a high-$\Rzero$ distinguished limit, and finite-$K$ stochastic-entry fluctuations. These are validity limits of the asymptotic regime rather than missing pieces of the main derivation.
\end{abstract}

\section{Executive formula sheet: predictions, definitions, and validity conditions}
\label{sec:executive}

This section is a self-contained statement of the result.  It is intended for a
reader who wants to implement or assess the approximation without following the
matched-asymptotic derivation.  The 2024 Parsons--Bolker--Dushoff--Earn burnout
paper carries the phase-plane matching through the refined term that is called
\emph{first order} here (the constant $C_0$); it does not contain the
$D_\theta$ second-order extension derived in this note.

\subsection{Basic inputs and quantities derived from them}

The independent dimensionless inputs are
\[
 \Rzero>1,\qquad 0<\rho\ll1,\qquad 0\leq\theta\leq1,
 \qquad K\gg1.
\]
Here $\Rzero$ is the basic reproduction number, $\rho$ is the susceptible
recruitment rate divided by the infectious removal rate, $\theta$ specifies the
recruitment law, and $K$ is population size.  Fractions of susceptible and
infectious hosts are denoted by $x$ and $y$.  Define, in the following order,
\begin{align*}
 h_\theta(x)&=(1-x)x^\theta,
 &F(x)&=1-x+\frac{1}{\Rzero}\log x,\\
 \xs&=\frac1{\Rzero},
 &\xf&=-\frac1{\Rzero}W_0(-\Rzero e^{-\Rzero}).
\end{align*}
Thus $\xf\in(0,\xs)$ is the nontrivial root of $F(\xf)=0$.  The endemic
infectious fraction and the standard matching height are
\begin{equation}
 y^*=\rho h_\theta(\xs)
 =\rho(\Rzero-1)\Rzero^{-(1+\theta)},
 \qquad
 \BL=\sqrt{\frac{y^*}{K}}.
 \label{eq:exec-ybl}
\end{equation}
The alternative choice $\BL=y^*$ used in some comparison plots may be inserted
in every formula below, but it is an empirical fallback: because it is
$O(\rho)$ rather than $o(\rho)$, it does not have the same formal second-order
error guarantee as~\eqref{eq:exec-ybl}.

The refined first-order constant is obtained entirely from $(\Rzero,\theta)$:
\begin{align}
 \widetilde a_\theta&=\frac{h_\theta(\xf)}{\xf},\nonumber\\
 \Jfun&=\int_{\xf}^{1}\left[
 \frac{(\xs-u)h_\theta(u)}{u^2F(u)}
 -\frac{\widetilde a_\theta}{u-\xf}\right]du,\nonumber\\
 \Cfun&=\frac{(1-\xf)(\xs-\xf)}{\xf}
 \exp\!\left(\frac{\Jfun}{\widetilde a_\theta}\right).
 \label{eq:exec-C}
\end{align}
The integrand in $J_\theta$ has a removable endpoint singularity.

For the second-order constant, define the coefficient integrands
\begin{align*}
 q_1(u)&=-\frac{h_\theta(u)(\Rzero u-1)}
 {\Rzero^2u^2F(u)},\\
 q_2(u)&=\frac{h_\theta(u)(\Rzero u-1)}
 {\Rzero^2u^2F(u)^2}
 \left[Y_1(u)-\frac{h_\theta(u)}{\Rzero u}\right].
\end{align*}
Then the two outer coefficients are the nested one-dimensional integrals
\begin{equation}
 \boxed{\qquad
 Y_1(x)=\int_1^x q_1(u)\,du,
 \qquad
 Y_2(x)=\int_1^x q_2(u)\,du.
 \qquad}
 \label{eq:exec-Y12}
\end{equation}
These are exactly equivalent to $Y_i'=q_i$, $Y_i(1)=0$.  Numerically, an ODE
integrator is simply an adaptive quadrature engine for these cumulative
integrals.  The apparent singularity at $u=1$ is removable.  Near $u=\xf$ the
raw second integral is badly conditioned; the production calculation therefore
subtracts its known singular part, integrates the finite remainder, and adds
back the analytic primitive, as detailed in Section~\ref{sec:D-numerics}.
For $0<y<F(\xs)$ let $X_0(y)\in(\xf,\xs)$ be the descending-branch solution
of $F(X_0)=y$, and set
\begin{align*}
 X_1(y)&=-\frac{Y_1(X_0)}{F'(X_0)},\\
 X_2(y)&=-\frac{\tfrac12F''(X_0)X_1^2+Y_1'(X_0)X_1+Y_2(X_0)}
 {F'(X_0)}.
\end{align*}
Define the local constants
\begin{align*}
 \delta&=1-\Rzero\xf,&h_f&=h_\theta(\xf),
 &h_f'&=\xf^{\theta-1}\{\theta-(\theta+1)\xf\},\\
 a&=\frac{h_f}{\delta},&b&=\frac{\Rzero\xf}{\delta},
 &c&=\delta h_f'+\Rzero h_f,\\
 \kappa&=\frac{h_fc}{2\delta^3},
 &f_1&=\frac{\delta}{\Rzero\xf},
 &f_2&=-\frac1{\Rzero\xf^2},\\
 \alpha&=\frac{h_f}{\Rzero\xf},
 &\beta&=\alpha\log\frac{f_1}{C_\theta},
 &p_0&=\frac{2\delta\xf h_f'-(1+2\delta)h_f}
 {2\Rzero\delta\xf^2},\\
 A_{-1}&=\frac{h_f/\xf-c}{\delta^2},
 &B_0&=\beta-\alpha,
 &B_1&=-\frac{h_f}{2\Rzero\delta\xf^2}.
\end{align*}
For $s>0$, subtract the explicitly known singular part
\begin{equation}
 q_{2,\mathrm{sing}}(s)
 =-\frac{a\alpha\log s+aB_0}{s^2}
 +\frac{A_{-1}\alpha\log s+A_{-1}B_0-aB_1}{s},
 \label{eq:exec-q2sing}
\end{equation}
and define
\begin{align}
 q_{2,\mathrm{reg}}(u)&=q_2(u)-q_{2,\mathrm{sing}}(u-\xf),\nonumber\\
 P_2(s)&=\frac{a\alpha\log s+a\beta}{s}
 +\frac{A_{-1}\alpha}{2}\log^2s
 +(A_{-1}B_0-aB_1)\log s,\nonumber\\
 K_2&=-P_2(1-\xf)-\int_{\xf}^{1}q_{2,\mathrm{reg}}(u)\,du,\nonumber\\
 N_{20}&=K_2+\frac{\alpha\beta f_2}{f_1^2}
 -\frac{(\alpha+\beta)p_0}{f_1}
 +\frac{f_2\beta^2}{2f_1^2}.
 \label{eq:exec-D-parts}
\end{align}
The production formula for the second-order constant is therefore
\begin{equation}
 \boxed{\qquad
 D_\theta=-\frac{N_{20}}{f_1}-\kappa\frac{\beta^2}{\alpha^2}.
 \qquad}
 \label{eq:exec-D}
\end{equation}
The integrand $q_{2,\mathrm{reg}}(u)=O(\log(u-\xf))$ is integrable at
$u=\xf$, so~\eqref{eq:exec-D} requires one finite regularized quadrature and
no small-$y$ extrapolation.  For general $\theta$ this is the explicit
semi-analytical expression; an elementary closed form is not available.  The
implementation stops a short distance from each endpoint and supplies the
omitted pieces from their local log--power expansions.  It also normalizes the
numerically integrated $Y_1$ so that its finite part is exactly
$\beta=\alpha\log(f_1/C_\theta)$.  These are quadrature-stabilization steps
only: the endpoint distances, fitted local coefficients, and normalization
shift cancel from the limiting quantity and are not extra parameters in
$D_\theta$.
The
equivalent limit characterization
\begin{equation}
 D_\theta=\lim_{y\downarrow0}\left[
 X_2(y)-\kappa\log^2\!\left(\frac{C_\theta}{y}\right)\right]
 \label{eq:exec-D-limit}
\end{equation}
is useful for the derivation and cross-checking only, not for production
evaluation.  The cancellation proving equivalence is given in
Section~\ref{sec:D-numerics}.

\subsection{First- and second-order entry points}

Only differences of the action are needed.  Define them directly by
\begin{equation}
 \Afun(x)-\Afun(z)=\int_z^x
 \frac{1-\Rzero u}{(1-u)u^\theta}\,du,
 \qquad \xf\leq z,x\leq\xs,
 \label{eq:exec-action}
\end{equation}
and put
\begin{align*}
 L_{\rm BL}&=\log(C_\theta/\BL),
 &Q&=\frac{\Rzero a-bc}{a\delta^2},\\
 M_1&=\rho L_{\rm BL}+\frac ba\BL,\\
 M_2&=M_1+\rho^2\frac{D_\theta}{a}
 +\rho\BL Q(L_{\rm BL}+1)+\frac{bQ}{2a}\BL^2.
\end{align*}
The direct corner expansions for the entry point are, explicitly,
\begin{equation}
 \boxed{
 x_{\rm in}^{(1),\,\mathrm{dir}}
 =\xf+\rho aL_{\rm BL}+b\BL,}
 \label{eq:exec-xin-direct1}
\end{equation}
and
\begin{equation}
 \boxed{\begin{aligned}
 x_{\rm in}^{(2),\,\mathrm{dir}}
 ={}&\xf+\rho aL_{\rm BL}+b\BL
 +\rho^2\{\kappa L_{\rm BL}^2+D_\theta\}\\
 &+\frac{\rho\BL}{\delta^2}
 \{\Rzero a(L_{\rm BL}+1)-bc\}
 +\frac{b\Rzero}{2\delta^2}\BL^2.
 \end{aligned}}
 \label{eq:exec-xin-direct2}
\end{equation}
These two formulas display the asymptotic corrections transparently.  For
production calculations we recommend the action-resummed versions: they keep
the same first- or second-order matching information but avoid unnecessarily
Taylor-expanding the nonlinear action.  The recommended first- and second-order
entry estimates are therefore the unique admissible roots
\begin{equation}
 \boxed{\quad
 \Afun(x_{\rm in}^{(j)})-\Afun(\xf)=M_j,
 \qquad \xf<x_{\rm in}^{(j)}<\xs,
 \quad j=1,2.\quad}
 \label{eq:exec-xin}
\end{equation}
They exist precisely when
\begin{equation}
 0<M_j<\Delta A_f,
 \qquad \Delta A_f=\Afun(\xs)-\Afun(\xf).
 \label{eq:exec-admissible}
\end{equation}
Failure of~\eqref{eq:exec-admissible} is a validity diagnostic, not a license
to extrapolate the scalar equation beyond its matching interval.

\subsection{Burnout and persistence probabilities}

For either $x_{\rm in}=x_{\rm in}^{(1)}$ or $x_{\rm in}^{(2)}$, define the exact
one-dimensional Kendall integral
\begin{equation}
 \mathcal I(x_{\rm in})=\int_{x_{\rm in}}^1
 \frac{1}{\rho(1-X)X^\theta}
 \exp\!\left\{\frac{\Afun(X)-\Afun(x_{\rm in})}{\rho}\right\}dX.
 \label{eq:exec-Kendall}
\end{equation}
The production calculation evaluates the same integral in the overflow-safe
form
\begin{align}
 \Lambda(x_{\rm in})
 &=\frac{\Afun(\xs)-\Afun(x_{\rm in})}{\rho},\nonumber\\
 J(x_{\rm in})
 &=\int_{x_{\rm in}}^1
 \frac{1}{\rho(1-X)X^\theta}
 \exp\!\left\{\frac{\Afun(X)-\Afun(\xs)}{\rho}\right\}dX,\nonumber\\
 \log\mathcal I&=\Lambda+\log J.
 \label{eq:exec-Kendall-scaled}
\end{align}
Thus $J$ is only a rescaled numerical integral, not an additional
approximation or an unreported model quantity.
With $m=K\BL$ infectives at entry, the conditional probabilities are
\begin{equation}
 \boxed{\quad P_{\rm burn\mid est}\simeq
 \left(\frac{\mathcal I}{1+\mathcal I}\right)^{K\BL},
 \qquad
 P_{\rm persist\mid est}\simeq1-P_{\rm burn\mid est}.\quad}
 \label{eq:exec-prob}
\end{equation}
Numerically these are evaluated without forming a potentially enormous
$\mathcal I$:
\begin{equation}
 B=m\log\!\left(1+e^{-\log\mathcal I}\right),
 \qquad
 P_{\rm persist\mid est}=-\operatorname{expm1}(-B),
 \qquad P_{\rm burn\mid est}=e^{-B}.
 \label{eq:exec-prob-stable}
\end{equation}
For one initially infectious individual, the early establishment probability
is $p_{\rm est}\simeq1-1/\Rzero$, so the unconditional first-trough persistence
probability is
\begin{equation}
 P_{\rm persist,uncond}\simeq
 \left(1-\frac1{\Rzero}\right)P_{\rm persist\mid est}.
 \label{eq:exec-uncond}
\end{equation}
For $I_0$ independent initial lineages, replace the leading factor by
$1-\Rzero^{-I_0}$.

The optional Laplace closed form uses
\begin{equation}
 \Lambda_j=\frac{\Delta A_f-M_j}{\rho},
 \qquad
 \mathcal I_{L,j}=\sqrt{\frac{2\pi}{\Rzero y^*}}e^{\Lambda_j},
 \label{eq:exec-Laplace}
\end{equation}
followed by~\eqref{eq:exec-prob} with $\mathcal I_{L,j}$ in place of
$\mathcal I$.  The further simplification
$-\log P_{\rm burn\mid est}\sim K\BL/\mathcal I_{L,j}$ additionally requires
$\mathcal I_{L,j}\gg1$.

\paragraph{Which route requires $x_{\rm in}$?}
There are two distinct computational pipelines.
\begin{enumerate}
 \item \textbf{Exact Kendall (recommended high-accuracy route):}
 compute $C_\theta,D_\theta$ and $M_j$; solve
 $\Afun(x_{\rm in}^{(j)})-\Afun(\xf)=M_j$ for $x_{\rm in}^{(j)}$; then use
 that value as the lower limit in~\eqref{eq:exec-Kendall}.  Thus this route
 genuinely requires $x_{\rm in}$.
 \item \textbf{Laplace closed form:} compute the action directly as
 \[
 \Lambda_j=\frac{\Delta A_f-M_j}{\rho}.
 \]
 Then use~\eqref{eq:exec-Laplace} and~\eqref{eq:exec-prob}.  No root solve for
 $x_{\rm in}$ and no separate evaluation of $\Afun(x_{\rm in})$ are required,
 because the matching equation has already supplied their action difference.
\end{enumerate}

\subsection{Conditions: what is required, and why}

The conditions are best treated as diagnostics rather than as a single vague
statement that every parameter must be ``small'' or ``large.''

\paragraph{Common to both orders.}
\begin{enumerate}
 \item $\rho\ll1$ with fixed $\Rzero>1$, and in practice
 $(\Rzero-1)/\rho\gg1$.  This separates the fast epidemic wave from slow
 susceptible replenishment and avoids the nonuniform near-threshold limit.
 \item $1/K\ll\BL\ll\rho$ (hence $K\BL\gg1$ and $Ky^*\gg1$ for the standard
 choice).  The first inequality makes deterministic entry concentrated enough
 to match to a branching process; the second makes infection feedback on the
 susceptible path negligible in the Kendall layer.
 \item The deterministic descending trajectory must actually cross
 $y=\BL$, and~\eqref{eq:exec-admissible} must hold.  Otherwise there is no
 entry point of the type assumed by the calculation.
 \item The fixed-$\Rzero$ expansion must not be pushed into the joint
 high-$\Rzero$ regime.  For $0\leq\theta<1$, a useful warning test is
 $\xf\gg\rho^{1/(1-\theta)}$; if these scales compete, replenishment changes
 the final-size corner itself.  The endpoint $\theta=1$ has a different
 distinguished balance and must not be judged by this formula.
\end{enumerate}

\paragraph{First-order matching.}
In addition to the common conditions, the corner displacement must be small:
\begin{equation}
 \rho|L_{\rm BL}|+\BL\ll \xs-\xf,
 \qquad 0<M_1<\Delta A_f.
 \label{eq:exec-first-cond}
\end{equation}
This is the operational requirement that the logarithmic switchback and the
finite-height correction remain local to the final-size corner.  It is the
appropriate accuracy level corresponding to the refined formula in the
original burnout paper when $\theta=0$.

\paragraph{Second-order matching.}
All first-order conditions remain necessary.  In addition, the nominal
second-order correction
\begin{equation}
 \Delta M_2=M_2-M_1
 =\rho^2\frac{D_\theta}{a}+\rho\BL Q(L_{\rm BL}+1)
 +\frac{bQ}{2a}\BL^2
\end{equation}
must be hierarchically smaller than the leading displacement, for example
\begin{equation}
 |\Delta M_2|\ll |M_1|,
 \qquad 0<M_2<\Delta A_f.
 \label{eq:exec-second-cond}
\end{equation}
This check explains why a formal second-order curve can be worse than the
first-order curve at moderate $\rho$: the extra terms are then no longer small
corrections and may over-correct.  Second order improves deterministic matching
only; it cannot cure finite-$K$ stochastic-entry error or a poor Laplace
approximation.

\paragraph{Laplace approximation only.}
The exact Kendall integral~\eqref{eq:exec-Kendall} needs no saddle-point
assumption.  Replacing it by~\eqref{eq:exec-Laplace} additionally requires
$\Lambda_j\gg1$, so that the integrand is sharply localized near $x=\xs$.
When $\Lambda_j=O(1)$, retain the exact integral even if the entry estimate is
second order.

\section{The problem: from a deterministic epidemic tail to stochastic burnout}
We consider the dimensionless SIR-type dynamics in a single patch
\begin{equation}
\label{eq:model}
\frac{dx}{d\sigma}=\rho h_\theta(x)-\Rzero xy,
\qquad
\frac{dy}{d\sigma}=(\Rzero x-1)y,
\end{equation}
where
\begin{equation}
\label{eq:h}
h_\theta(x)=(1-x)x^\theta,
\qquad 0\leq\theta\leq1.
\end{equation}
Here $x$ and $y$ are the susceptible and infectious fractions, respectively, $\Rzero>1$ is the basic reproduction number, and $\rho$ controls the slow time scale of susceptible replenishment. The case $\theta=0$ gives the standard linear replenishment model, while $\theta=1$ gives logistic replenishment; intermediate values interpolate continuously between the two.

At low susceptible density,
\[
h_\theta(x)\sim x^\theta,
\qquad x\to0,
\]
so $\theta$ may be interpreted as the low-density elasticity of recruitment with respect to the current susceptible fraction. The range $0\leq\theta\leq1$ therefore includes both the Parsons model and the logistic limit.

Define the epidemic growth threshold
\begin{equation}
\label{eq:xstar}
\xs=\frac1{\Rzero}.
\end{equation}
The endemic equilibrium satisfies $x=\xs$ and
\begin{equation}
\label{eq:ystar}
y^*=\rho h_\theta(\xs)
=\rho(\Rzero-1)\Rzero^{-(1+\theta)}.
\end{equation}

Let $K$ denote the characteristic host population size of one patch. After the first epidemic wave, the number of infectious individuals falls from a macroscopic level into a stochastic branching regime. To match the deterministic epidemic tail to the low-count stochastic process, we use the practical matching height
\begin{equation}
\label{eq:ybl}
\BL=\sqrt{\frac{y^*}{K}}.
\end{equation}
When $Ky^*\gg1$,
\[
\frac1K\ll\BL\ll y^*\sim\rho,
\]
so $\BL$ lies in an overlap region between the one-infectious-individual scale and the deterministic endemic prevalence.

We define $\xin$ as the susceptible fraction at the first downward crossing of $y=\BL$ after the first epidemic peak. Our goal is to construct the map
\[
(\Rzero,\rho,\theta,K)
\longrightarrow \xin
\longrightarrow \Pone,
\]
without integrating the full two-dimensional system~\eqref{eq:model} for every parameter set.

\section{The Parsons framework and the logic of the present extension}
Parsons and Earn constructed matched asymptotic approximations for the standard SIR model with vital dynamics, including the epidemic outer solution, the left corner, the $x$-axis boundary layer, and an additional boundary layer required in a joint high-$\Rzero$ limit\cite{PE2024}. Parsons et al. then connected this phase-plane construction to a time-inhomogeneous branching process to obtain the probability of burnout during the first epidemic trough\cite{PNAS2024}.

The present work uses the same singular-perturbation architecture but relies only on two structural facts:
\begin{enumerate}
\item when $\rho=0$, the geometry of the first epidemic wave is independent of $h_\theta$;
\item when $y$ is small, susceptible replenishment enters the slow basin dynamics through $h_\theta(x)$.
\end{enumerate}
Hence the same outer/corner hierarchy can be extended from standard linear replenishment to the whole family $h_\theta(x)=(1-x)x^\theta$.

The main asymptotic regime considered here is
\begin{equation}
\label{eq:limit}
\rho\to0,
\qquad \Rzero>1\ \text{fixed},
\qquad K\ \text{large enough that }1/K\ll\BL\ll\rho.
\end{equation}
The joint high-$\Rzero$ limit and finite-$K$ stochastic entry are treated separately as boundaries of validity.

\section{Leading order: universal epidemic geometry}
Set $\rho=0$. Equation~\eqref{eq:model} gives
\[
\frac{dy}{dx}=\frac1{\Rzero x}-1.
\]
Starting from a naive invasion $(x,y)=(1,0^+)$, the first-wave outer trajectory is
\begin{equation}
\label{eq:F}
F(x)=1-x+\xs\log x,
\qquad y=F(x).
\end{equation}
The epidemic peak occurs at $x=\xs$. The zero-demography final susceptible fraction $\xf$ is the smaller root of $F(\xf)=0$:
\begin{equation}
\label{eq:xf}
\xf=-\frac1{\Rzero}W_0(-\Rzero e^{-\Rzero}),
\end{equation}
where $W_0$ is the principal branch of the Lambert $W$ function.

The point $\xf$ is the anchor of the corner matching, but it cannot be used directly as the finite-$\rho$ value of $\xin$. The reason is that the later burnout probability depends on an action scaled by $1/\rho$; even if $\xin-\xf\to0$, that difference may still be exponentially amplified. A systematic matched correction to $\xin$ is therefore required.

\section{Exact phase-plane equation and the direct perturbation hierarchy}
From~\eqref{eq:model}, the exact phase-plane equation is
\begin{equation}
\label{eq:phase}
\frac{dy}{dx}
=\frac{(\Rzero x-1)y}{\rho h_\theta(x)-\Rzero xy}.
\end{equation}
In the first-wave outer region, write the direct expansion
\begin{equation}
\label{eq:Yexp}
y=F(x)+\rho Y_1(x)+\rho^2Y_2(x)+O(\rho^3).
\end{equation}
Substituting~\eqref{eq:Yexp} into~\eqref{eq:phase} and equating successive powers of $\rho$ gives
\begin{equation}
\label{eq:Y1}
Y_1'(x)
=-\frac{h_\theta(x)(\Rzero x-1)}{\Rzero^2x^2F(x)},
\qquad Y_1(1)=0,
\end{equation}
and
\begin{equation}
\label{eq:Y2}
Y_2'(x)
=\frac{h_\theta(x)(\Rzero x-1)}{\Rzero^2x^2F(x)^2}
\left[Y_1(x)-\frac{h_\theta(x)}{\Rzero x}\right],
\qquad Y_2(1)=0.
\end{equation}
The apparent singularities at $x=1$ are removable. Up to second order, the outer correction is therefore determined by two one-dimensional coefficient equations.

An equivalent bookkeeping device uses the first integral of the $\rho=0$ SIR system,
\[
H=x+y-\frac1{\Rzero}\log x.
\]
Along the full model,
\[
\frac{dH}{d\sigma}
=\rho h_\theta(x)\left(1-\frac1{\Rzero x}\right).
\]
If one writes $H=1+\rho H_1+\rho^2H_2+\cdots$, then $H_1=Y_1$ and $H_2=Y_2$. Thus the conserved-quantity formulation is simply an equivalent organization of the same direct phase-plane hierarchy, not a separate approximation route.

\section{First-order matching: the logarithmic singularity and $C_\theta$}
\subsection{Relation to the Parsons notation}
Parsons and Earn use two visually similar small parameters.  Their first
definition, introduced in the time-parametrized SIR equations, is
\[
\varepsilon=\frac{\mu}{\gamma+\mu},
\]
the infectious period divided by the host lifetime.  In the standard
recruitment case $\theta=0$, their equation is
\[
 \frac{dx}{d\sigma}=\varepsilon(1-x)-\Rzero xy,
\]
whereas ours is $dx/d\sigma=\rho(1-x)-\Rzero xy$.  Therefore the exact
dictionary is
\begin{equation}
 \boxed{\varepsilon=\rho.}
 \label{eq:eps-vareps-map}
\end{equation}
They then introduce the second glyph
\begin{equation}
 \boxed{\epsilon=\frac{\varepsilon}{\Rzero}
 =\frac{\rho}{\Rzero},}
 \label{eq:eps-map}
\end{equation}
because it simplifies the phase-plane equations.  Thus $\varepsilon$ and
$\epsilon$ are distinct, with $\varepsilon=\Rzero\epsilon$; neither is an
additional biological parameter.  This note expands directly in $\rho$,
\[
 y=F+\rho Y_1+\rho^2Y_2+\cdots,
\]
while Parsons and Earn organize the corresponding phase-plane expansion in
$\epsilon$.  Since $\Rzero$ is fixed, $\epsilon=O(\rho)$ and the asymptotic
orders agree.  To translate their first-order coefficient to ours, set
\[
 G_\theta=\Rzero Y_1.
\]
Then
\begin{equation}
\label{eq:Gprime}
G_\theta'(x)
=\frac{(\xs-x)h_\theta(x)}{x^2F(x)}.
\end{equation}
Define
\begin{equation}
\label{eq:atilde}
\widetilde a_\theta
=\frac{h_\theta(\xf)}{\xf}
=(1-\xf)\xf^{\theta-1}.
\end{equation}
Because
\[
F'(\xf)=\frac{\xs-\xf}{\xf},
\]
we have, as $x\to\xf^+$,
\[
G_\theta'(x)\sim\frac{\widetilde a_\theta}{x-\xf}.
\]
Thus the first-order outer correction necessarily contains a logarithmic singularity.

\subsection{The regularized finite part}
Define the regularized integral
\begin{equation}
\label{eq:Jtheta}
\Jfun
=\int_{\xf}^{1}
\left[
\frac{(\xs-u)h_\theta(u)}{u^2F(u)}
-\frac{\widetilde a_\theta}{u-\xf}
\right]du.
\end{equation}
The two singular terms cancel at $u=\xf$, so~\eqref{eq:Jtheta} is a smooth one-dimensional integral. It yields the refined first-order matching amplitude
\begin{equation}
\label{eq:Ctheta}
\Cfun
=\frac{(1-\xf)(\xs-\xf)}{\xf}
\exp\!\left(\frac{\Jfun}{\widetilde a_\theta}\right).
\end{equation}
The constant $C_\theta$ depends only on $(\Rzero,\theta)$ and is independent of $(\rho,K)$.

When $\theta=0$,~\eqref{eq:Jtheta} reduces to the regularized integral in the refined Parsons et al. burnout estimate, and~\eqref{eq:Ctheta} reduces to its first-order prefactor\cite{PNAS2024}. Hence $C_\theta$ is a direct generalization of the standard burnout matching constant.

\section{Inverse outer expansion and corner scaling}
Let $X_0(y)$ denote the descending branch satisfying $F(X_0)=y$, and write
\begin{equation}
\label{eq:Xexp}
X(y;\rho)=X_0(y)+\rho X_1(y)+\rho^2X_2(y)+O(\rho^3).
\end{equation}
Implicit-function expansion gives
\begin{equation}
\label{eq:X1}
X_1(y)=-\frac{Y_1(X_0)}{F'(X_0)},
\end{equation}
and
\begin{equation}
\label{eq:X2}
X_2(y)
=-\frac{\frac12F''(X_0)X_1^2+Y_1'(X_0)X_1+Y_2(X_0)}{F'(X_0)}.
\end{equation}

When $y=O(\rho)$, the raw outer series is no longer uniform. Introduce the corner scaling
\begin{equation}
\label{eq:corner_scale}
y=\rho v,
\qquad
x=\xf+\rho u_1(v)+\rho^2u_2(v)+\cdots.
\end{equation}
Define
\begin{equation}
\label{eq:abcd}
\delta=1-\Rzero\xf,
\qquad
h_f=h_\theta(\xf),
\qquad
h_f'=h_\theta'(\xf),
\end{equation}
and
\begin{equation}
\label{eq:abc2}
a=\frac{h_f}{\delta},
\qquad
b=\frac{\Rzero\xf}{\delta},
\qquad
c=\delta h_f'+\Rzero h_f,
\qquad
\kappa=\frac{h_fc}{2\delta^3}.
\end{equation}
For $h_\theta(x)=(1-x)x^\theta$,
\[
h_\theta'(x)=x^{\theta-1}\{\theta-(\theta+1)x\}.
\]

Let
\begin{equation}
\label{eq:L}
L=\log\frac{C_\theta}{y}.
\end{equation}
Expanding the exact phase-plane equation~\eqref{eq:phase} directly in the corner gives, at first order,
\begin{equation}
\label{eq:u1}
u_1(v)=aL+bv.
\end{equation}
The term $\rho aL$ is the logarithmic switchback, while $by$ is the finite-height correction. Equation~\eqref{eq:u1} is exactly the corner representation of the refined first-order matching obtained above.

\section{Second-order matching: $D_\theta$ and the next switchback}
\label{sec:D-numerics}
Continuing the same corner expansion to $O(\rho^2)$ gives
\begin{equation}
\label{eq:u2}
\begin{split}
u_2(v)={}&\kappa L^2+D_\theta
+\frac{\Rzero a}{\delta^2}(L+1)v
-\frac{bc}{\delta^2}v
+\frac{b\Rzero}{2\delta^2}v^2.
\end{split}
\end{equation}
The $L^2$ term is the natural next-order continuation of the first-order logarithmic switchback.

Both $C_\theta$ and the new second-order matching constant $D_\theta$ are uniquely determined by regularized limits of the inverse outer coefficients:
\begin{equation}
\label{eq:Climit}
\log C_\theta
=\lim_{y\downarrow0}
\left[\frac{X_1(y)}a+\log y\right],
\end{equation}
\begin{equation}
\label{eq:Dlimit}
D_\theta
=\lim_{y\downarrow0}
\left[
X_2(y)-\kappa\log^2\!\left(\frac{C_\theta}{y}\right)
\right].
\end{equation}
Equation~\eqref{eq:Climit} gives the same $C_\theta$ as the explicit regularized integral~\eqref{eq:Ctheta}; this is an important internal consistency check between the two representations. The constant $D_\theta$ contains the genuinely new second-order information. It can be computed without a small-$y$ extrapolation, as follows.

Put $s=x-\xf$, $f_1=F'(\xf)=\delta/(\Rzero\xf)$,
$f_2=F''(\xf)=-1/(\Rzero\xf^2)$, and
\begin{equation}
\label{eq:Y1local}
Y_1(x)=\alpha\log s+\beta+p_0s+O(s^2),
\quad
\alpha=\frac{h_f}{\Rzero\xf},
\quad
\beta=\alpha\log\frac{f_1}{C_\theta},
\end{equation}
where
\[
p_0=\frac{2\delta\xf h_f'-(1+2\delta)h_f}
{2\Rzero\delta\xf^2}.
\]
Define
\[
A_{-1}=\frac{h_f/\xf-c}{\delta^2},
\qquad B_0=\beta-\alpha,
\qquad B_1=-\frac{h_f}{2\Rzero\delta\xf^2}.
\]
The integrand on the right of~\eqref{eq:Y2}, denoted by $q_2(x)$, has singular part
\begin{equation}
\label{eq:q2sing}
q_{2,{\rm sing}}(s)
=-\frac{a\alpha\log s+aB_0}{s^2}
+\frac{A_{-1}\alpha\log s+A_{-1}B_0-aB_1}{s}.
\end{equation}
Direct expansion of the exact integrand gives
$q_{2,{\rm reg}}(x):=q_2(x)-q_{2,{\rm sing}}(x-\xf)=O(\log(x-\xf))$,
so the regularized integrand is integrable at $\xf$.  A convenient singular
primitive is
\begin{equation}
\label{eq:P2}
P_2(s)=\frac{a\alpha\log s+a\beta}{s}
+\frac{A_{-1}\alpha}{2}\log^2s
+(A_{-1}B_0-aB_1)\log s,
\end{equation}
for which direct differentiation gives $P_2'=q_{2,{\rm sing}}$.  Hence
\begin{equation}
\label{eq:K2}
K_2=-P_2(1-\xf)-\int_{\xf}^{1}q_{2,{\rm reg}}(u)\,du,
\qquad
Y_2(x)=P_2(x-\xf)+K_2+o(1).
\end{equation}

For completeness, write $T=\alpha\log s+\beta$.  From~\eqref{eq:X1},
\[
X_1=-\frac{T}{f_1}
+s\left(-\frac{p_0}{f_1}+\frac{f_2T}{f_1^2}\right)+O(s^2\log s).
\]
Thus the singular part of $Y_1'X_1$ is
$-\alpha T/(f_1s)=-(a\alpha\log s+a\beta)/s$, which cancels both the
$(\log s)/s$ and $1/s$ terms in $Y_2$.  The remaining finite numerator in
~\eqref{eq:X2} has constant part
\begin{equation}
\label{eq:N20}
N_{20}=K_2+\frac{\alpha\beta f_2}{f_1^2}
-\frac{(\alpha+\beta)p_0}{f_1}
+\frac{f_2\beta^2}{2f_1^2}.
\end{equation}
Its logarithmic terms reduce to
$-f_1\kappa(\log s+\beta/\alpha)^2$; since
$\log(C_\theta/y)=-\log s-\beta/\alpha+o(1)$, the logarithmic part of $X_2$
is exactly $\kappa\log^2(C_\theta/y)$.  The finite-quadrature formula is therefore
\begin{equation}
\label{eq:Dquadrature}
D_\theta=-\frac{N_{20}}{f_1}-\kappa\frac{\beta^2}{\alpha^2}.
\end{equation}
Equations~\eqref{eq:K2} and~\eqref{eq:Dquadrature} are the recommended numerical
definition: they require one finite regularized quadrature and prove finiteness
directly.  Equation~\eqref{eq:Dlimit} remains the equivalent asymptotic
characterization, but no endpoint extrapolation is needed in the production
calculation.

\section{Action formulation: from the corner expansion to a stable scalar equation}
\subsection{The action primitive}
Define
\begin{equation}
\label{eq:Adef}
\Afun'(x)
=\frac{1-\Rzero x}{h_\theta(x)}
=\frac{1-\Rzero x}{(1-x)x^\theta}.
\end{equation}
In the $x$-axis basin, where $y\ll\rho$, infection-induced depletion of susceptible individuals is of higher order, so
\[
\frac{dx}{d\sigma}\simeq\rho h_\theta(x),
\qquad
\frac{d\log y}{dx}
=-\frac1\rho\Afun'(x).
\]
Thus $\Afun$ is the natural variable linking deterministic epidemic matching to the later branching-process exponent.

For $\theta\neq1$, one may take
\begin{equation}
\label{eq:Ahyper}
\Afun(x)
=\frac{x^{1-\theta}}{1-\theta}
\left[
\Rzero-(\Rzero-1)
{}_2F_1(1,1-\theta;2-\theta;x)
\right].
\end{equation}
Important special cases are
\begin{align}
\mathcal A_0(x)
&=\Rzero x+(\Rzero-1)\log(1-x),\\
\mathcal A_{1/2}(x)
&=2\Rzero\sqrt{x}-2(\Rzero-1)\operatorname{artanh}(\sqrt{x}),\\
\mathcal A_1(x)
&=\log x+(\Rzero-1)\log(1-x).
\end{align}

\subsection{Refined first-order action equation}
Rather than truncating $x$ directly, we reinsert~\eqref{eq:u1} into the full action. This gives the practical refined first-order equation
\begin{equation}
\label{eq:M1}
\Afun(x_{\rm in}^{(1)})-\Afun(\xf)
=\rho L_{\rm BL}+\frac ba\BL,
\qquad
L_{\rm BL}=\log\frac{C_\theta}{\BL}.
\end{equation}
Because
\[
\frac ba=\frac{\Rzero\xf}{h_f},
\]
this is exactly the finite-$\BL$ first-order formula.

\subsection{Second-order action-resummed master equation}
Define
\begin{equation}
\label{eq:Q}
Q=\frac{\Rzero a-bc}{a\delta^2}.
\end{equation}
Substitute $x=\xf+\rho u_1+\rho^2u_2$ into $\Afun(x)$ and retain all terms required to second order under the corner scaling. A key simplification is that the $L^2$ term in $u_2$ cancels the quadratic contribution from the Taylor expansion of $\Afun$. The result is
\begin{equation}
\label{eq:M2}
\begin{split}
\Afun(x_{\rm in}^{(2)})-\Afun(\xf)
={}&\rho L_{\rm BL}
+\frac ba\BL
+\rho^2\frac{D_\theta}{a}\\
&+\rho\BL Q(L_{\rm BL}+1)
+\frac{bQ}{2a}\BL^2.
\end{split}
\end{equation}
This is the recommended second-order master equation.

Equation~\eqref{eq:M2} has two important properties. First, it is a systematic higher-order extension of~\eqref{eq:M1}:
\[
M_2=M_1+\text{second-order corrections}.
\]
Second, it remains only a scalar root-finding problem for $x_{\rm in}^{(2)}$. For fixed $(\Rzero,\theta)$, the constants $C_\theta$ and $D_\theta$ need to be precomputed only once.

It is also important that the practical cutoff satisfies $\BL\ll\rho$ but does not generally satisfy $\BL\ll\rho^2$. Therefore, if one wants uniform second-order accuracy at finite $K$, it is not consistent to retain only the pure $\rho^2D_\theta$ term while discarding the $\rho\BL$ or $\BL^2$ terms.

\section{The Kendall layer: from $x_{\rm in}$ to burnout probability}
Once the trajectory enters the low-prevalence layer, treat the susceptible path as a deterministic background and neglect the feedback of the rare infection on $x$:
\[
\frac{dx}{d\sigma}\simeq\rho h_\theta(x).
\]
Each infectious lineage then has birth rate $\Rzero x(\sigma)$ and removal rate $1$. Kendall's result for a time-inhomogeneous linear birth--death process gives the extinction probability of one lineage as
\begin{equation}
\label{eq:q}
q_1=\frac{\mathcal I}{1+\mathcal I},
\end{equation}
where
\begin{equation}
\label{eq:Kendall}
\mathcal I(\xin)
=\int_{\xin}^{1}
\frac1{\rho h_\theta(X)}
\exp\!\left[
\frac{\Afun(X)-\Afun(\xin)}{\rho}
\right]dX.
\end{equation}
If the matching line contains
\[
m=K\BL
\]
infectious individuals, then the conditional probability of burnout is
\begin{equation}
\label{eq:Pburn}
\Pburn
\simeq\left(\frac{\mathcal I}{1+\mathcal I}\right)^{K\BL},
\end{equation}
and the probability of persistence through the first trough is
\begin{equation}
\label{eq:P1}
\Pone\simeq1-\Pburn.
\end{equation}
A numerically more stable quantity is the burnout exponent
\begin{equation}
\label{eq:B}
B=-\log\Pburn
=K\BL\log\left(1+\frac1{\mathcal I}\right).
\end{equation}

Hence the high-accuracy semi-analytical pipeline is
\[
(\Rzero,\theta)\to(\xf,C_\theta,D_\theta),
\quad
(\rho,K)\to\BL,
\quad
\eqref{eq:M2}\to x_{\rm in}^{(2)},
\quad
\eqref{eq:Kendall}\to B,\Pone.
\]
The entire procedure requires only one-dimensional coefficient equations, a scalar root solve, and one one-dimensional Kendall integral; it does not require integrating the full two-dimensional epidemic ODE for each parameter set.

\section{The Laplace layer: a fully asymptotic closed form}
Define the action
\begin{equation}
\label{eq:Lambda}
\Lambda(\xin)
=\frac{\Afun(\xs)-\Afun(\xin)}{\rho}.
\end{equation}
The stationary point of the exponent in the Kendall integrand is determined by
\[
1-\Rzero x=0,
\]
so for every $\theta$ the unique internal saddle lies at
\[
\xs=\frac1{\Rzero}.
\]
Applying Laplace's method around $\xs$ gives
\begin{equation}
\label{eq:Ilap}
\mathcal I
\sim
\sqrt{\frac{2\pi}{\Rzero g(\xs)}}e^\Lambda,
\qquad
g(\xs)=\rho h_\theta(\xs)=y^*.
\end{equation}
Substituting this approximation into~\eqref{eq:B} gives the finite-$m$ Laplace version; if in addition $\mathcal I\gg1$, then
\begin{equation}
\label{eq:Blap}
B
\sim
K\BL
\sqrt{\frac{\Rzero g(\xs)}{2\pi}}
e^{-\Lambda}.
\end{equation}

The action-resummed formulation has a further computational advantage. Let $M_2$ denote the right-hand side of~\eqref{eq:M2}, and define
\[
\Delta A_f=\Afun(\xs)-\Afun(\xf).
\]
Then one need not first solve for $x_{\rm in}^{(2)}$ in order to obtain
\begin{equation}
\label{eq:LambdaM2}
\Lambda^{(2)}
=\frac{\Delta A_f-M_2}{\rho}.
\end{equation}
Thus the fully asymptotic prediction for $P_1$ can be computed directly from the matching constants and the action. The root solve for $x_{\rm in}^{(2)}$ is required only when one wants the exact Kendall integral.

\section{Why both exact Kendall and Laplace versions should be retained}
Increasing the asymptotic order of $x_{\rm in}$ controls only the deterministic matching error; it cannot repair the saddle-point approximation itself. When
\[
\Lambda\gg1,
\]
the Kendall integrand is sharply concentrated near $x_*$ and the Laplace formula is accurate. When $\Lambda=O(1)$ or smaller, the Laplace approximation may have substantial error even if $x_{\rm in}$ is already extremely accurate.

The two versions should therefore be viewed as different levels of approximation:
\begin{enumerate}
\item \textbf{Primary numerical/semi-analytical prediction}: second-order $x_{\rm in}$ plus the exact one-dimensional Kendall integral;
\item \textbf{Fully asymptotic closed form}: second-order action plus Laplace's method, valid when the action is sufficiently large.
\end{enumerate}
Numerical validation should report ``matching error'' and ``Laplace-only error'' separately.

\section{Domain of validity and identifiable failure modes}
The present theory addresses deterministic-entry asymptotics for fixed $\Rzero>1$, small $\rho$, and large $K$. Outside that regime, several distinct failure mechanisms can occur.

\subsection{No boundary-layer entry}
If the first deterministic trough satisfies
\[
y_{\min}>\BL,
\]
the descending branch of the first epidemic wave never reaches the matching line, so the present definition of $x_{\rm in}$ does not apply. Such parameter values should be labeled \emph{no boundary-layer entry}, rather than assigned a large approximation error.

\subsection{$\Rzero\to1^+$}
When $\Rzero$ is very close to one, the first epidemic wave becomes weak and the scales associated with $x_f$, $x_*$, the peak, and the trough begin to merge. The fixed-$\Rzero$ outer/corner hierarchy is then no longer uniform; a near-threshold distinguished limit is required.

\subsection{Small-action Laplace breakdown}
Even with an accurate $x_{\rm in}$, Laplace's approximation may fail if $\Lambda$ is not large. In that regime the correct response is to retain the exact Kendall integral, rather than to increase the formal order of the $x_{\rm in}$ expansion.

\subsection{Joint high-$\Rzero$ limit}
If $\Rzero$ also becomes large so that $x_f$ competes with $\rho$ or smaller scales, a new boundary layer is required. Since
\[
h_\theta(x)\sim x^\theta,
\qquad x\to0,
\]
the local balance between replenishment and infection-induced consumption,
\[
\rho x^\theta\sim x,
\]
gives the natural scale
\begin{equation}
\label{eq:highRscale}
x\sim\rho^{1/(1-\theta)},
\end{equation}
for $0\leq\theta<1$. The case $\theta=1$ corresponds to a different distinguished limit. This is a future high-$\Rzero$ uniformization problem, not simply an omitted ordinary $O(\rho^3)$ term in~\eqref{eq:M2}.

\subsection{Finite-$K$ stochastic entry}
The present derivation treats $x_{\rm in}$ as deterministic. For finite $K$, the actual entry point is a random variable $X_{\rm in}$, and its fluctuations may be exponentially amplified because
\[
\Lambda'(x)
=-\frac{1-\Rzero x}{\rho h_\theta(x)}.
\]
This effect is distinct from deterministic matching error and should be assessed separately through the distribution of entry points in stochastic simulations.

\section{Theoretical chain and computational interface}
At this stage the large-population deterministic-entry problem forms a closed chain:
\[
(\Rzero,\rho,\theta,K)
\longrightarrow
(\xf,C_\theta,D_\theta,\BL)
\longrightarrow
x_{\rm in}^{(1)},x_{\rm in}^{(2)}
\longrightarrow
\mathcal I,\Lambda
\longrightarrow
\Pburn,\Pone.
\]

A practical implementation naturally separates into a precomputation layer and a parameter layer:
\begin{enumerate}
\item For each $(\Rzero,\theta)$, compute $\xf$, $C_\theta$, and $D_\theta$, using the regularized quadrature~\eqref{eq:Dquadrature} for the latter;
\item for each $(\rho,K)$, compute $y^*$ and $\BL$, and use~\eqref{eq:M1} and~\eqref{eq:M2} to obtain first- and second-order predictions;
\item for the high-accuracy version, use the exact Kendall integral~\eqref{eq:Kendall};
\item when a closed form is needed, use the action~\eqref{eq:LambdaM2} and the Laplace approximation~\eqref{eq:Ilap};
\item record no-entry status, $\Lambda$, and the difference between first and second order as validity diagnostics.
\end{enumerate}

There are no undefined intermediate quantities left in the theory. The next stage is systematic numerical validation: compare first order, second order, the deterministic ODE benchmark, exact Kendall, Laplace, and full stochastic simulations, and use these comparisons to map the domain of validity quantitatively.

\appendix
\section{Two equivalent representations of the first-order constant $C_\theta$}
The first-order matching constant can be obtained explicitly from the singular finite part of $G_\theta$ near $x_f$, which yields~\eqref{eq:Ctheta}, or from the inverse outer coefficient through~\eqref{eq:Climit}. The two are equivalent because
\[
Y_1=\frac{G_\theta}{\Rzero}
\]
and $X_1$ is the implicit-function inversion of the same first-order correction. This equivalence is useful computationally: the explicit regularized integral gives a stable first-order calculation, while the inverse-coefficient representation connects seamlessly to the second-order $D_\theta$ hierarchy.

\section{Core diagnostics to retain in numerical validation}
To separate the main sources of approximation error, the numerical study should retain at least
\[
\frac{y_{\min}}{\BL},
\qquad
\frac{x_{\rm in}^{(1)}-x_{\rm in}^{\rm ODE}}{x_{\rm in}^{\rm ODE}},
\qquad
\frac{x_{\rm in}^{(2)}-x_{\rm in}^{\rm ODE}}{x_{\rm in}^{\rm ODE}},
\]
\[
\Delta\Lambda_1,
\qquad
\Delta\Lambda_2,
\qquad
\Lambda_{\rm ODE},
\qquad
B_{\rm Kendall},
\qquad
B_{\rm Laplace}.
\]
The regularized quadrature for $D_\theta$ should also be tested for endpoint-treatment and numerical-precision sensitivity. The constant $C_\theta$ should be cross-validated using both~\eqref{eq:Ctheta} and~\eqref{eq:Climit}.

\begin{thebibliography}{9}
\bibitem{PE2024}
T. L. Parsons and D. J. D. Earn,
\emph{Uniform Asymptotic Approximations for the Phase Plane Trajectories of the SIR Model with Vital Dynamics},
SIAM Journal on Applied Mathematics 84(4), 1580--1608 (2024).

\bibitem{PNAS2024}
T. L. Parsons, B. M. Bolker, J. Dushoff, and D. J. D. Earn,
\emph{The probability of epidemic burnout in the stochastic SIR model with vital dynamics},
Proceedings of the National Academy of Sciences 121(5), e2313708120 (2024).
\end{thebibliography}

\end{document}
